% Part: methods
% Chapter: proofs
% Section: inference-patterns

\documentclass[../../../include/open-logic-section]{subfiles}

\begin{document}

\olfileid[hi]{mth}{prf}{pat}

\olsection{अनुमान प्रतिरूप}

प्रमाण अलग-अलग अनुमानों से बने होते हैं। अनुमान करते समय हम सामान्यतः
``अतः'', ``इस प्रकार'' या ``फलतः'' जैसे शब्द से उसका संकेत करते हैं। अनुमान
प्रायः प्रमाण में पहले से उपलब्ध एक या दो तथ्यों पर निर्भर करता है---वह कुछ
ऐसा हो सकता है जिसे हमने माना हो, अथवा पहले ही किसी अनुमान से निष्कर्षित
किया हो। स्पष्टता के लिए हम इन बातों पर चिह्न लगा सकते हैं और अनुमान में
बताते हैं कि किन अन्य कथनों का उपयोग कर रहे हैं। अनुमान में प्रायः यह
व्याख्या भी होगी कि नया निष्कर्ष पिछली बातों से \emph{क्यों} निष्पन्न होता
है। अनुमान के कुछ सामान्य प्रतिरूप प्रमाणों में बहुत बार प्रयुक्त होते हैं;
नीचे उनमें से कुछ देखेंगे। आगमन द्वारा प्रमाण जैसे कुछ अनुमान प्रतिरूप अधिक
विस्तृत हैं (और उनकी चर्चा बाद में होगी)।

हम एक अनुमान प्रतिरूप पर पहले ही चर्चा कर चुके हैं: किसी परिभाषा को खोलना
या लागू करना। परिभाषा खोलते समय हम परिभाषक का उपयोग करके परिभाष्य वाली बात
को फिर कहते हैं। उदाहरणतः मान लें हमने प्रमाण के दौरान पहले ही $D=E$ (a)
स्थापित किया है। तब समुच्चयों के लिए $=$ की परिभाषा लागू करके अनुमान कर सकते
हैं: ``अतः परिभाषा से (a) के अनुसार $D$ का प्रत्येक !!{element}, $E$ का
!!a{element} है और इसके विपरीत भी।''

कुछ भ्रमकारी रूप से हम अनुमान का औचित्य प्रायः अनुमान करते समय नहीं, बल्कि
उससे पहले लिखते हैं। मान लें हमने $D=E$ अभी सिद्ध नहीं किया, किंतु करना
चाहते हैं। यदि लक्ष्य निष्कर्ष $D=E$ है, तो परिभाषा लागू करके इस लक्ष्य को
फिर कह सकते हैं: $D=E$ सिद्ध करने के लिए सिद्ध करना होगा कि $D$ का प्रत्येक
!!{element}, $E$ का !!a{element} है और इसके विपरीत। अतः प्रमाण का रूप होगा:
(a) सिद्ध करें कि $D$ का प्रत्येक !!{element}, $E$ का !!a{element} है;
(b) $E$ का प्रत्येक !!{element}, $D$ का !!a{element} है; (c) इसलिए (a) और
(b) से $=$ की परिभाषा द्वारा $D=E$। किंतु सामान्यतः हम इसे इस ढंग से नहीं
लिखेंगे। इसके बदले ऐसा कुछ लिख सकते हैं,
\begin{quote}
हम $D=E$ दिखाना चाहते हैं। $=$ की परिभाषा से इसका अर्थ यह दिखाना है कि
$D$ का प्रत्येक !!{element}, $E$ का !!a{element} है और इसके विपरीत।

(a) \dots (यह प्रमाण कि $D$ का प्रत्येक !!{element}, $E$ का
!!a{element} है) \dots

(b) \dots (यह प्रमाण कि $E$ का प्रत्येक !!{element}, $D$ का
!!a{element} है) \dots
\end{quote}

\subsection{संयोजन का उपयोग}

शायद सरलतम अनुमान प्रतिरूप संयोजन के किसी एक संयोज्य को निष्कर्ष के रूप में
निकालना है। दूसरे शब्दों में: यदि हमने $p$ और $q$ माना अथवा पहले ही सिद्ध
किया है, तो $p$ (और $q$ भी) अनुमानित कर सकते हैं। यह इतना मूल अनुमान है कि
प्रायः इसका उल्लेख नहीं होता। उदाहरणतः $D=E$ की परिभाषा खोलने पर हमने
स्थापित किया कि $D$ का प्रत्येक !!{element}, $E$ का !!a{element} है और इसके
विपरीत। इससे निष्कर्षित कर सकते हैं कि $E$ का प्रत्येक !!{element}, $D$ का
!!a{element} है (यही ``इसके विपरीत'' वाला भाग है)।

\subsection{संयोजन सिद्ध करना}

कभी-कभी सिद्ध करने को दी गई बात संयोजन रूप की होगी; आपसे ``$p$ और $q$
सिद्ध करें'' कहा जाएगा। इस स्थिति में केवल दो काम करने हैं: $p$ सिद्ध करें,
फिर $q$ सिद्ध करें। प्रमाण को दो अनुभागों में बाँटकर स्पष्टता के लिए उन पर
चिह्न लगा सकते हैं। पहली टिप्पणियाँ बनाते समय पृष्ठ के ऊपर ``(1) $p$ सिद्ध
करें'' और पृष्ठ के मध्य में ``(2) $q$ सिद्ध करें'' लिख सकते हैं। (निस्संदेह
हो सकता है आपसे स्पष्ट रूप से संयोजन सिद्ध करने को न कहा जाए, किंतु प्रमाण
में संयोजन सिद्ध करना आवश्यक मिले। उदाहरणतः यदि $D=E$ सिद्ध करने को कहा जाए,
तो $=$ की परिभाषा खोलने के बाद सिद्ध करना होगा: $D$ का प्रत्येक
!!{element}, $E$ का !!a{element} है \emph{और} $E$ का प्रत्येक !!{element},
$D$ का !!a{element} है।)

\subsection{वियोजन सिद्ध करना}

जब सिद्ध की जा रही बात वियोजन रूप की हो (अर्थात् ``$p$ अथवा $q$'' रूप का
कथन हो), तो किसी एक वियोज्य को सत्य दिखाना पर्याप्त है। किंतु लगभग कभी ऐसा
नहीं होता कि कोई एक वियोज्य आपके प्रमेय की मान्यताओं से सीधे निष्पन्न हो।
अधिकतर प्रमेय की मान्यताएँ स्वयं वियोजक होती हैं, या आप दिखा रहे होते हैं कि
किसी प्रकार की सभी वस्तुओं में दो में से एक गुण है, किंतु कुछ में पहला और
अन्य में दूसरा गुण है। यहाँ स्थितियों द्वारा प्रमाण उपयोगी होता है (नीचे
देखें)।

\subsection{सशर्त प्रमाण}

आपको मिलने वाले अनेक प्रमेय सशर्त रूप में होंगे (अर्थात् दिखाएँ कि यदि $p$
सत्य है, तो $q$ भी सत्य है)। इन स्थितियों को व्यवस्थित करना सरल है---सशर्त
का पूर्ववर्ती (यहाँ $p$) मान लें और उससे निष्कर्ष $q$ सिद्ध करें। अतः यदि
प्रमेय कहता है, ``यदि $p$, तो $q$'', प्रमाण ``$p$ मानें'' से आरंभ करें और
अंत तक $q$ सिद्ध हो जाना चाहिए।

सशर्तों को भिन्न ढंग से कहा जा सकता है। अतः ``यदि $p$, तो $q$'' के बदले
प्रमेय कह सकता है ``$p$ केवल यदि $q$'', ``$q$ यदि $p$'', या ``$q$, बशर्ते
$p$''। इन सबका अर्थ समान है और इनमें $p$ मानकर उस मान्यता से $q$ सिद्ध करना
होता है। याद करें कि द्विसशर्त (``$p$ तभी और केवल तभी जब $q$'') वास्तव में
दो सशर्तों का मेल है: यदि $p$, तो $q$, और यदि $q$, तो $p$। तब केवल सशर्त
प्रमाण के दो दृष्टांत करने हैं: एक पहले सशर्त के लिए और दूसरा दूसरे के लिए।
किंतु कभी-कभी कई अन्य ``तभी और केवल तभी'' कथनों को जोड़कर ऐसा कथन सिद्ध किया
जा सकता है, जिससे $p$ से आरंभ कर $q$ पर समाप्त हों---पर उस स्थिति में
सुनिश्चित करना होगा कि प्रत्येक चरण वास्तव में ``तभी और केवल तभी'' है।

\subsection{सार्वत्रिक दावे}

सार्वत्रिक दावे का उपयोग सरल है: यदि कोई बात हर वस्तु के लिए सत्य है, तो
वह प्रत्येक विशिष्ट वस्तु के लिए सत्य है। अतः यदि प्रमाण की परिकल्पना
$A\subseteq B$ है, तो इसका अर्थ ($\subseteq$ की परिभाषा खोलकर) है कि
प्रत्येक $x\in A$ के लिए $x\in B$। इसलिए यदि पहले से जानते हैं कि $z\in A$,
तो $z\in B$ निष्कर्षित कर सकते हैं।

सार्वत्रिक दावा सिद्ध करना कुछ कठिन लग सकता है। सामान्यतः ये कथन इस रूप के
होते हैं: ``यदि $x$ में $P$ है, तो उसमें $Q$ है'' या ``सभी $P$, $Q$ हैं।''
निस्संदेह कथन इस रूप में पूरी तरह न बैठे, और ठीक क्या सिद्ध करने को कहा गया
है यह जानने में कुछ अभ्यास लगता है। किंतु प्रायः हमें सिद्ध करना होता है कि
किसी गुण वाली सभी वस्तुओं में कोई अन्य विशिष्ट गुण है।

सार्वत्रिक दावा सिद्ध करने का ढंग उस एक गुण वाली वस्तुओं के लिए नाम या चर
प्रस्तुत करना और फिर दिखाना है कि उनमें दूसरा गुण भी है। इसे ऐसे कह सकते हैं
कि \emph{सभी} $P$ के लिए कोई बात सिद्ध करने हेतु उसे किसी \emph{मनमाने} $P$
के लिए सिद्ध करना होगा। और प्रस्तुत नाम किसी मनमाने $P$ का नाम है। मनमानी
वस्तुओं के इन नामों के लिए सामान्यतः एक अक्षर उपयोग करते हैं और अक्षर प्रायः
परंपराओं का पालन करते हैं: जैसे प्राकृतिक संख्याओं के लिए $n$, !!{formula}ों
के लिए $!A$, समुच्चयों के लिए $A$, फलनों के लिए $f$, इत्यादि।

युक्ति पूरे प्रमाण में सामान्यता बनाए रखना है। आरंभ में मानें कि किसी मनमानी
वस्तु (``$x$'') में गुण $P$ है, और केवल परिभाषाओं या अनुमत मान्यताओं के
आधार पर दिखाएँ कि $x$ में गुण $Q$ है। चूँकि यह निर्धारित नहीं किया कि $x$
विशेषतः क्या है, सिवाय इसके कि उसमें $P$ है, तब कह सकते हैं कि $P$ वाली हर
वस्तु में $Q$ है। संक्षेप में $x$, गुण $P$ वाली \emph{सभी} वस्तुओं का
प्रतिनिधि है।

\begin{prop}
  सभी समुच्चयों $A$ और $B$ के लिए $A\subseteq A\cup B$।
\end{prop}

\begin{proof}
  $A$ और $B$ मनमाने समुच्चय हों। हम $A\subseteq A\cup B$ दिखाना चाहते
  हैं। $\subseteq$ की परिभाषा से इसका अर्थ है: प्रत्येक $x$ के लिए यदि
  $x\in A$, तो $x\in A\cup B$। अतः $x\in A$, $A$ का कोई मनमाना
  !!{element} हो। हमें दिखाना है कि $x\in A\cup B$। चूँकि $x\in A$,
  $x\in A$ अथवा $x\in B$। अतः
  $x\in\Setabs{x}{x\in A\lor x\in B}$। किंतु $\cup$ की परिभाषा से इसका
  अर्थ $x\in A\cup B$ है।
\end{proof}


\subsection{स्थितियों द्वारा प्रमाण}

मान लें मान्यता या पहले से स्थापित निष्कर्ष के रूप में आपके पास वियोजन
है---आपने माना या सिद्ध किया है कि $p$ अथवा $q$ सत्य है। आप $r$ सिद्ध करना
चाहते हैं। यह दो चरणों में करें: पहले $p$ को सत्य मानकर $r$ सिद्ध करें, फिर
$q$ को सत्य मानकर $r$ फिर सिद्ध करें। यह इसलिए काम करता है कि हम दोनों
विकल्पों में से एक को सत्य मानते या जानते हैं। दोनों चरण स्थापित करते हैं
कि दोनों में से कोई भी $r$ की सत्यता के लिए पर्याप्त है। (यदि दोनों सत्य
हैं, तो $r$ के सत्य होने के एक नहीं दो कारण हैं। $p$ और $q$ दोनों मानकर $r$
को अलग से सिद्ध करना आवश्यक नहीं।) अपने काम का संकेत देने के लिए हम घोषणा
करते हैं कि ``हम स्थितियों में भेद करते हैं।'' उदाहरणतः मान लें हम जानते हैं
कि $x\in B\cup C$। $B\cup C$ को
$\Setabs{x}{x\in B\text{ अथवा }x\in C}$ परिभाषित किया गया है। दूसरे शब्दों
में परिभाषा से $x\in B$ अथवा $x\in C$। इससे $x\in A$ सिद्ध करने के लिए पहले
$x\in B$ मानकर इस मान्यता से $x\in A$ सिद्ध करेंगे, फिर $x\in C$ मानकर उससे
$x\in A$ फिर सिद्ध करेंगे। मान्यता के नीचे ``हम स्थितियों में भेद करते हैं'',
फिर उसके नीचे ``स्थिति (1): $x\in B$'', और पृष्ठ के बीच में ``स्थिति (2):
$x\in C$'' लिखेंगे। फिर पृष्ठ के ऊपर और नीचे के आधे भाग भरेंगे।

स्थितियों द्वारा प्रमाण विशेषतः तब उपयोगी है जब सिद्ध की जा रही बात स्वयं
वियोजक हो। एक सरल उदाहरण देखें:

\begin{prop}
मान लें $B\subseteq D$ और $C\subseteq E$। तब
$B\cup C\subseteq D\cup E$।
\end{prop}

\begin{proof}
  (a) $B\subseteq D$ और (b) $C\subseteq E$ मानें। परिभाषा से कोई भी
  $x\in B$, $D$ में भी है (c), और कोई भी $x\in C$, $E$ में भी है (d)।
  $B\cup C\subseteq D\cup E$ दिखाने के लिए ($\subseteq$ की परिभाषा से)
  दिखाना होगा कि यदि $x\in B\cup C$, तो $x\in D\cup E$। $x\in B\cup C$
  तभी जब $x\in B$ अथवा $x\in C$ ($\cup$ की परिभाषा से)। इसी प्रकार
  $x\in D\cup E$ तभी जब $x\in D$ अथवा $x\in E$। अतः हमें दिखाना है: किसी
  भी $x$ के लिए यदि $x\in B$ अथवा $x\in C$, तो $x\in D$ अथवा $x\in E$।

  \begin{quote}
  अब तक हमने केवल परिभाषाएँ खोली हैं!{} हमने अपनी प्रतिज्ञप्ति को
  $\subseteq$ और $\cup$ के बिना फिर सूत्रबद्ध किया और अब एक सार्वत्रिक
  सशर्त दावा सिद्ध करना शेष है। ऊपर की चर्चा के अनुसार यह ऐसा $x$ मानकर
  किया जाता है जिसके बारे में ``यदि'' वाला भाग सत्य मानते हैं, और फिर
  दिखाते हैं कि ``तो'' वाला भाग भी सत्य है। दूसरे शब्दों में हम मानेंगे कि
  $x\in B$ अथवा $x\in C$ और दिखाएँगे कि $x\in D$ अथवा $x\in E$।
  \footnote{यह अनुच्छेद केवल हमारे काम की व्याख्या करता है---यह प्रमाण का
  भाग नहीं और अपना प्रमाण लिखते समय आपको इतने विस्तार में जाने की आवश्यकता
  नहीं।}
  \end{quote}

  मान लें $x\in B$ अथवा $x\in C$। हमें दिखाना है कि $x\in D$ अथवा
  $x\in E$। हम स्थितियों में भेद करते हैं।

  स्थिति 1: $x\in B$। (c) से $x\in D$। अतः $x\in D$ अथवा $x\in E$।
  (यहाँ हमने पिछले उपखंड में चर्चित अनुमान किया है!)
      
  स्थिति 2: $x\in C$। (d) से $x\in E$। अतः $x\in D$ अथवा $x\in E$।
 \end{proof}


\subsection{अस्तित्व दावा सिद्ध करना}

अस्तित्व दावा सिद्ध करने को कहे जाने पर प्रश्न सामान्यतः ``सिद्ध करें कि
कोई $x$ है ऐसा कि $\dots x\dots$'' रूप का होगा, अर्थात् कोई वस्तु है जिसमें
``$\dots x\dots$'' से वर्णित गुण है। इस स्थिति में किसी उपयुक्त वस्तु की
पहचान कर दिखाना होगा कि उसमें आवश्यक गुण है। यह सीधा लगता है, किंतु इस
प्रकार का प्रमाण कठिन हो सकता है। सामान्यतः इसमें किसी वस्तु का
\emph{निर्माण} या उसे \emph{परिभाषित} करना और सिद्ध करना शामिल है कि इस
प्रकार परिभाषित वस्तु में आवश्यक गुण है। सही वस्तु पाना कठिन हो सकता है,
उसमें आवश्यक गुण सिद्ध करना कठिन हो सकता है, और कभी-कभी यह दिखाना भी कठिन
होता है कि आपने कोई वस्तु परिभाषित करने में सफलता पाई है!{}

सामान्यतः आप वस्तु को निर्दिष्ट करके इसे लिखेंगे, जैसे ``$x$, \dots हो''
(जहाँ \dots{} मन की वस्तु निर्दिष्ट करता है), संभवतः यह सिद्ध करेंगे कि
$\dots$ वास्तव में किसी विद्यमान वस्तु का वर्णन करता है, और फिर दिखाएँगे कि
$x$ में गुण $Q$ है। एक सरल उदाहरण देखें।

\begin{prop}
  मान लें $x\in B$। तब कोई $A$ है ऐसा कि $A\subseteq B$ और
  $A\neq\emptyset$।
\end{prop}

\begin{proof}
  $x\in B$ मानें। $A=\{x\}$ रखें।
  \begin{quote}
    यहाँ हमने समुच्चय $A$ को उसके !!{element} गिनाकर परिभाषित किया है।
    चूँकि हम मानते हैं कि $x$ एक वस्तु है और !!{element} गिनाकर सदैव
    समुच्चय बना सकते हैं, हमें यहाँ यह दिखाने की आवश्यकता नहीं कि समुच्चय
    $A$ परिभाषित करने में सफल हुए हैं। किंतु अभी भी दिखाना होगा कि $A$ में
    प्रतिज्ञप्ति द्वारा अपेक्षित गुण हैं। इसके बिना प्रमाण पूर्ण नहीं!{}
  \end{quote}
  चूँकि $x\in A$, $A\neq\emptyset$।
  \begin{quote}
    यह $A$ की $\{x\}$ के रूप में परिभाषा और स्पष्ट तथ्यों
    $x\in\{x\}$ तथा $x\notin\emptyset$ पर निर्भर है।
  \end{quote}
  चूँकि $x$, $\{x\}$ का एकमात्र !!{element} है और $x\in B$, $A$ का
  प्रत्येक !!{element}, $B$ का भी !!a{element} है। $\subseteq$ की
  परिभाषा से $A\subseteq B$।
\end{proof}

\subsection{अस्तित्व दावों का उपयोग}

मान लें आप जानते हैं कि कोई अस्तित्व दावा सत्य है (आपने उसे सिद्ध किया है
या वह उपयोग योग्य परिकल्पना है), जैसे ``किसी $x$ के लिए $x\in A$'' अथवा
``कोई $x\in A$ है।'' प्रमाण में इसका उपयोग करने के लिए आप मान सकते हैं कि
आपके पास उन वस्तुओं में से एक का नाम है जिनका अस्तित्व आपकी परिकल्पना कहती
है। चूँकि $A$ में कम-से-कम एक वस्तु है, ऐसी वस्तुएँ हैं जिनके लिए वह नाम हो
सकता है। निस्संदेह आप किसी एक को चुन या उसका आगे वर्णन नहीं कर सकें (सिवाय
इसके कि वह $\in A$ है)। किंतु प्रमाण के उद्देश्य से मान सकते हैं कि आपने
उसे चुना और नाम दिया है। ऐसा नाम चुनना महत्त्वपूर्ण है जिसका पहले उपयोग न
किया हो (या जो परिकल्पनाओं में न आता हो), अन्यथा त्रुटि हो सकती है। प्रमाण
में इसका संकेत ``किसी $x$ के लिए $x\in A$'' से ``$a\in A$ हो'' पर जाकर
करते हैं। अब $a$ के बारे में तर्क कर सकते हैं, अन्य परिकल्पनाएँ उपयोग कर
सकते हैं, इत्यादि, जब तक किसी निष्कर्ष $p$ पर न पहुँचें। यदि $p$ में अब $a$
का उल्लेख नहीं, तो $p$ मान्यता $a\in A$ से स्वतंत्र है और आपने दिखाया कि
वह केवल ``किसी $x$ के लिए $x\in A$'' से निष्पन्न होता है।

\begin{prop}
यदि $A\neq\emptyset$, तो $A\cup B\neq\emptyset$।
\end{prop}

\begin{proof}
  मान लें $A\neq\emptyset$। अतः किसी $x$ के लिए $x\in A$।
  \begin{quote}
    यहाँ पहले हमने प्रतिज्ञप्ति की परिकल्पना को केवल फिर कहा। यह परिकल्पना,
    अर्थात् $A\neq\emptyset$, एक अस्तित्व दावा छिपाती है, जिस तक कुछ
    परिभाषाएँ खोलकर ही पहुँचते हैं। $=$ की परिभाषा बताती है कि
    $A=\emptyset$ तभी जब प्रत्येक $x\in A$, $\emptyset$ में भी हो और प्रत्येक
    $x\in\emptyset$, $A$ में भी हो। दोनों पक्षों का निषेध करने पर मिलता है:
    $A\neq\emptyset$ तभी जब या तो कोई $x\in A$, $\notin\emptyset$ हो, अथवा
    कोई $x\in\emptyset$, $\notin A$ हो। चूँकि कुछ भी $\in\emptyset$ नहीं,
    दूसरा वियोज्य कभी सत्य नहीं हो सकता, और ``$x\in A$ तथा
    $x\notin\emptyset$'' केवल $x\in A$ रह जाता है। अतः
    $x\neq\emptyset$ तभी जब किसी $x$ के लिए $x\in A$। यह अस्तित्व दावा है।
    अब $A$ के !!{element}ों में से एक के लिए नाम प्रस्तुत करके उस अस्तित्व
    दावे का उपयोग करते हैं:
  \end{quote}
  $a\in A$ हो।
  \begin{quote}
    अब हमने $\in A$ वस्तुओं में से एक के लिए नाम प्रस्तुत किया है। हम $a$
    के बारे में तर्क जारी रखेंगे, किंतु केवल $a\in A$ और कुछ नहीं मानने में
    सावधान रहेंगे:
  \end{quote}
  चूँकि $a\in A$, $\cup$ की परिभाषा से $a\in A\cup B$। अतः किसी $x$ के
  लिए $x\in A\cup B$, अर्थात् $A\cup B\neq\emptyset$।
  \begin{quote}
    अंतिम चरण में हम ``$a\in A\cup B$'' से ``किसी $x$ के लिए
    $x\in A\cup B$'' पर गए। इसमें अब $a$ का उल्लेख नहीं, अतः जानते हैं कि
    ``किसी $x$ के लिए $x\in A\cup B$'' केवल ``किसी $x$ के लिए $x\in A$''
    से निष्पन्न होता है। किंतु इसका अर्थ $A\cup B\neq\emptyset$ है।
  \end{quote}
\end{proof}

हमारी तरह ``$x$'' जैसे आबद्ध चरों को $a$ जैसे काल्पनिक नामों से अलग रखना
शायद अच्छी प्रथा है। किंतु व्यवहार में हम प्रायः ऐसा नहीं करते और केवल $x$
का उपयोग करते हैं, इस प्रकार:
\begin{quote}
मान लें $A\neq\emptyset$, अर्थात् कोई $x\in A$ है। $\cup$ की परिभाषा से
$x\in A\cup B$। अतः $A\cup B\neq\emptyset$।
\end{quote}
किंतु ऐसा करते समय विशेष सावधानी चाहिए कि भिन्न अस्तित्व दावों के लिए भिन्न
$x$ और $y$ उपयोग करें। उदाहरणतः निम्न ``यदि $A\neq\emptyset$ और
$B\neq\emptyset$, तो $A\cap B\neq\emptyset$'' (जो सत्य नहीं) का सही प्रमाण
\emph{नहीं} है।
\begin{quote}
मान लें $A\neq\emptyset$ और $B\neq\emptyset$। अतः किसी $x$ के लिए $x\in A$
और किसी $x$ के लिए $x\in B$ भी। चूँकि $x\in A$ और $x\in B$, $\cap$ की
परिभाषा से $x\in A\cap B$। अतः $A\cap B\neq\emptyset$।
\end{quote}
क्या आप पहचान सकते हैं कि गलत चरण कहाँ है और समझा सकते हैं कि परिणाम क्यों
सत्य नहीं है?

\end{document}
